\section{Conclusiones y Trabajos futuros}

La herramienta MDAA permite, de forma relativamente sencilla, implementar un flujo de trabajo MDA en el ámbito de la programación de sintetizadores de audio, proporcionando resultados muy interesantes al permitir la generación de código de forma automatizada y brindar trazabilidad entre sus distintos niveles de abstracción. Sin embargo, la herramienta está incompleta; solo brinda la posibilidad de que, tras la generación automática de una serie de ficheros estáticos, un modelo de razonamiento externo a ella sea capaz de generar el código fuente deseado. La evolución natural de esta herramienta debería ir enfocada a potenciar su autonomía, sin depender de modelos externos para la obtención de sus resultados finales e incluyendo en ella modelos locales entrenados para tal fin, que puedan encargarse también de la transformación automática entre los modelos DSL. MDAA parece brinda una base conceptual sólida para ello, y su evolución podría ser un interesante proyecto de fin de máster. Hay que tener en cuenta que MDAA es una herramienta en versión \textit{Alfa} todavía, por lo que quedaría mucho trabajo por realizar en ella para que pudiera ser considerada una herramienta robusta y mantenible. Las gramáticas DSL para el ámbito en testeo sí se pueden considerar maduras, en el sentido de que cubren la mayor parte de lo que un software de rango pequeño o mediano de este tipo suele requerir. Aun así, también se considera que pueden ser ampliamente mejorables, ya no en su estructura base, sino en las capacidades que ofrecen.

Recalcar también que, aunque el artículo se centra en un dominio específico, desarrollo de sintetizadores de audio software, la estructura de los DSL desarrollados, desde un punto de vista de alto nivel, enfocando esta mirada en su idea de definición de \textit{entes} (máquinas para el dominio ejemplo) en un DSL especifico, y sus interconexiones en otro distinto pero fuertemente relacionado, tanto para los niveles de abstracción de análisis como de diseño, parece un interesante punto de partida para su posible generalización, pudiendo partir de ese esquema inicial en el diseño de herramientas similares en otros ámbitos.

En definitiva, se considera este trabajo como un proyecto inicial, a completar, que podría abrir una línea de investigación interesante para futuros trabajos.

